\documentclass[11pt,a4paper]{article}

\usepackage[margin=2.3cm]{geometry}
\usepackage{amsmath,amssymb,amsthm,bm}
\usepackage{booktabs}
\usepackage{graphicx}
\usepackage{float}
\usepackage{listings}
\usepackage{xcolor}
\usepackage{caption}
\usepackage{microtype}
\usepackage{enumitem}
\usepackage[colorlinks=true,linkcolor=black,urlcolor=blue]{hyperref}

% ---------------------------------------------------------------- code style
\definecolor{codebg}{gray}{0.97}
\definecolor{codekw}{RGB}{0,90,160}
\definecolor{codecm}{RGB}{110,130,110}
\definecolor{codestr}{RGB}{160,70,40}

\lstdefinestyle{base}{
  backgroundcolor=\color{codebg},
  basicstyle=\ttfamily\footnotesize,
  keywordstyle=\color{codekw}\bfseries,
  commentstyle=\color{codecm}\itshape,
  stringstyle=\color{codestr},
  breaklines=true,
  breakatwhitespace=true,
  columns=fullflexible,
  showstringspaces=false,
  frame=single,
  rulecolor=\color{gray!40},
  framesep=4pt,
  xleftmargin=6pt,
  aboveskip=8pt,
  belowskip=8pt,
}
\lstdefinestyle{py}{style=base,language=Python}
\lstdefinelanguage{Stan}{
  morekeywords={data,parameters,model,generated,quantities,transformed,functions,
    int,real,vector,matrix,array,lower,upper,normal,bernoulli,bernoulli_logit,
    Phi,inv_logit,sqrt,target},
  morecomment=[l]{//},
  morestring=[b]",
}
\lstdefinestyle{stan}{style=base,language=Stan}

% ---------------------------------------------------------------- notation
\newcommand{\y}{\bm{y}}
\newcommand{\x}{\bm{x}}
\newcommand{\z}{\bm{z}}
\newcommand{\tv}{\bm{t}}
\newcommand{\Normal}{\mathrm{Normal}}
\newcommand{\Gam}{\mathrm{Gamma}}
\newcommand{\GP}{\mathrm{GP}}
\newcommand{\Bern}{\mathrm{Bernoulli}}
\newcommand{\E}{\mathbb{E}}
\newcommand{\Prob}{\mathbb{P}}
\newcommand{\ind}{\mathbf{1}}
\newcommand{\yrep}{\y^{\mathrm{rep}}}

% ---------------------------------------------------------------- code inclusion
% Segments are pulled verbatim from code/*.py between "# <<tag" and "# >>tag" markers,
% so the report can never drift out of step with the code that actually ran.
\lstset{rangeprefix=\#\ ,includerangemarker=false}
\newcommand{\pycode}[2]{\lstinputlisting[style=py,linerange=<<#2->>#2]{code/#1.py}}

% ---------------------------------------------------------------- generated results
% code/_report.py writes results/qN.tex as a list of \setresult{key}{value} lines.
% Cite a number with \result{q1.ppp}; unknown keys render as ?? rather than failing,
% so the document still compiles before the code has been run.
\makeatletter
\newcommand{\setresult}[2]{\expandafter\gdef\csname result@#1\endcsname{#2}}
\newcommand{\result}[1]{%
  \ifcsname result@#1\endcsname\csname result@#1\endcsname\else\textbf{??}\fi}
\makeatother
\newcommand{\loadresults}[1]{\IfFileExists{results/#1.tex}{\input{results/#1.tex}}{}}

% ---------------------------------------------------------------- marks
\newcommand{\partmarks}[1]{\hfill\normalfont\mbox{[#1 marks]}}
\newcommand{\questiontotal}[1]{%
  \vspace{6pt}\par\hfill\textbf{Question total: #1 Marks}\par}
\newcommand{\grandtotal}[1]{%
  \vspace{6pt}\par\hfill\textbf{Total: #1 Marks}\par}

\setlist[itemize]{itemsep=2pt,topsep=3pt}
\setlist[enumerate]{itemsep=2pt,topsep=3pt}

\title{\vspace{-1.4cm}\textbf{MATH70100 Bayesian Methods \& Computation}\\[4pt]
       \large Assessed Coursework 2}
\date{}

\loadresults{q1}
\loadresults{q2}
\loadresults{q3}

\begin{document}
\maketitle
\vspace{-1.6cm}

\noindent\rule{\textwidth}{0.4pt}
\vspace{2pt}

\noindent
TODO: preamble note --- state the computational environment (e.g.\ Python version, NumPy/SciPy/PyStan)
and name the accompanying notebooks from which the code segments below are reproduced. Record any
generative AI tools used as a learning aid, per the coursework's academic integrity policy.

\vspace{4pt}
\noindent\rule{\textwidth}{0.4pt}

% ==========================================================================
\section*{Question 1}
% ==========================================================================

\subsection*{(a) Bayes factors and posterior model probabilities \partmarks{3}}

Let
\begin{equation}
  p(\y \mid M) = \int p(\y \mid \beta,\sigma^2,M)\, p(\beta,\sigma^2\mid M)\,
  \mathrm d\beta\,\mathrm d\sigma^2
  \label{eq:marglik}
\end{equation}
denote the marginal likelihood of $\y$ under a model $M$. By
Definition~7.3.2 the \emph{Bayes factor} in favour of $M_i$ over $M_j$ is
\begin{equation}
B_{ij}(\y) \;=\; \frac{p(\y\mid M_i)}{p(\y\mid M_j)}
\;=\; \frac{\Prob(M_i\mid \y)}{\Prob(M_j\mid \y)}\Bigg/\frac{\Prob(M_i)}{\Prob(M_j)} .
\label{eq:bf}
\end{equation}

The second equality in \eqref{eq:bf} shows that $B_{ij}$ is exactly the factor by which the
data update the prior odds of $M_i$ against $M_j$ into the posterior odds.
\textcolor{red}{[TODO --- one sentence on why this is important. The prior model weights have
cancelled, so $B_{ij}$ depends on the data alone, therefore analysts who disagree on the priors still
agree on $B_{ij}$.]}
Values $B_{ij}>1$ mean the data favour $M_i$ over $M_j$; $B_{ij}<1$ favours $M_j$.Note
$B_{ij}=1/B_{ji}$ and $B_{ik}=B_{ij}B_{jk}$, so the three pairwise comparisons among
$M_1,M_2,M_3$ are determined by any two of them --- in practice one computes all evidences
once and forms whichever ratios are of interest.

\textcolor{red}{[TODO --- agaain explain why a conventional scale is needed at all. $B_{ij}$ is
unbounded above and carries no units, so unlike a probability it has no self-evident scale.]}
Their magnitude is therefore read off the conventional scale of Kass \& Raftery (1995),
reproduced in Table~7.1 of the notes.

\subsection*{(b) Monte Carlo estimation of the posterior predictive $p$-value \partmarks{4}}

TODO

\subsection*{(c) Marginal likelihoods, posterior model probabilities, preferred model \partmarks{5}}

TODO --- describe the computation, then the code follows.

\pycode{q1_model_choice}{q1c}

TODO --- report the three log marginal likelihoods and posterior probabilities. Numbers are
cited from the generated results file, e.g.\ $\log p(\y\mid M_1)=\result{q1.logml.M1}$ with
$\Prob(M_1\mid\y)=\result{q1.post.M1}$, so the text updates whenever the code is re-run. The
preferred model is $\result{q1.preferred}$.

\subsection*{(d) Fitting the preferred model and the posterior predictive check \partmarks{5}}

TODO

\pycode{q1_model_choice}{q1d}

TODO --- comment on adequacy. The estimated posterior predictive $p$-value is
$\hat p = \result{q1.ppp}$ (Monte Carlo standard error $\result{q1.ppp.se}$).

% Uncomment once code/q1_model_choice.py has written the figure.
% \begin{figure}[h]
% \centering
% \includegraphics[width=\textwidth]{figs/q1_ppc.pdf}
% \caption*{TODO: caption.}
% \end{figure}

\subsection*{(e) Most probable, yet poorly fitting \partmarks{3}}

TODO

\questiontotal{20}

\clearpage
% ==========================================================================
\section*{Question 2}
% ==========================================================================

\subsection*{(a) Posterior predictive distribution of $f(t^\star)$ \partmarks{5}}

TODO

\subsection*{(b) Interpretation of $\alpha$, $\rho$ and $\sigma^2$ \partmarks{3}}

TODO

\subsection*{(c) Posterior mean, credible bands and $\Prob(f(5)>0\mid\tv,\y)$ \partmarks{5}}

TODO --- describe the computation, then the code follows.

\pycode{q2_gp_regression}{q2c}

TODO --- report $\Prob(f(5)>0\mid\tv,\y)=\result{q2.prob.positive}$ and discuss.

% \begin{figure}[h]
% \centering
% \includegraphics[width=0.82\textwidth]{figs/q2_gp.pdf}
% \caption*{TODO: caption.}
% \end{figure}

\subsection*{(d) The Bayes action under asymmetric linear loss \partmarks{4}}

TODO --- identify the Bayes action as a posterior quantile, then evaluate it.

\pycode{q2_gp_regression}{q2d}

TODO --- the Bayes action is $d^\star = \result{q2.bayes.action}$.

\subsection*{(e) Effect of the length-scale: Model A ($\rho=1.5$) versus Model B ($\rho=0.4$) \partmarks{3}}

TODO

\questiontotal{20}

\clearpage
% ==========================================================================
\section*{Question 3}
% ==========================================================================

\subsection*{(a) Posterior density up to proportionality \partmarks{3}}

TODO

\subsection*{(b) Non-conjugacy, and why Stan is convenient \partmarks{2}}

TODO

\subsection*{(c) Stan model code \partmarks{6}}

TODO --- explain the model blocks, then the code follows.

\lstinputlisting[style=stan]{code/probit_regression.stan}

\subsection*{(d) Fitting the Stan model in Python \partmarks{3}}

TODO --- describe the fit and its convergence diagnostics.

\pycode{q3_probit_stan}{q3d}

\subsection*{(e) Posterior estimates \partmarks{3}}

\pycode{q3_probit_stan}{q3e}

TODO --- report $\E[\beta\mid\x,\z]=\result{q3.e.beta}$,
$\Prob(\beta>0\mid\x,\z)=\result{q3.p.beta.positive}$ and
$\Prob(z^\star=1\mid x^\star=0.75,\x,\z)=\result{q3.p.zstar}$.

\subsection*{(f) Probit versus logistic regression \partmarks{3}}

TODO

\questiontotal{20}

\vspace{4pt}
\grandtotal{60}

\end{document}
